% TEMPORARY PROOF APPENDIX. DO NOT \input THIS FILE INTO paper.tex.
% This file is compiled in isolation. It records algebraic proofs that must be
% reviewed before any part is moved into the manuscript.
\documentclass[11pt]{article}

\usepackage[margin=1in]{geometry}
\usepackage{amsmath,amssymb,amsthm,mathtools}
\usepackage{booktabs}
\usepackage{microtype}
\usepackage[nameinlink,noabbrev]{cleveref}

\newcommand{\F}{\mathbb F}
\newcommand{\A}{\mathcal A}
\newcommand{\Sym}{\operatorname{Sym}}
\newcommand{\Gr}{\operatorname{Gr}}
\newcommand{\rad}{\operatorname{rad}}
\newcommand{\im}{\operatorname{im}}
\newcommand{\Span}{\operatorname{Span}}
\newcommand{\codim}{\operatorname{codim}}
\newcommand{\Hom}{\operatorname{Hom}}

\theoremstyle{plain}
\newtheorem{theorem}{Theorem}[section]
\newtheorem{proposition}[theorem]{Proposition}
\newtheorem{lemma}[theorem]{Lemma}
\newtheorem{corollary}[theorem]{Corollary}
\theoremstyle{definition}
\newtheorem{definition}[theorem]{Definition}
\theoremstyle{remark}
\newtheorem{remark}[theorem]{Remark}

\title{Temporary proof appendix for the transverse and fixed map analyses}
\author{}
\date{12 August 2026}

\begin{document}
\maketitle

\begin{center}
\fbox{\parbox{0.9\textwidth}{
\centering
\textbf{DO NOT INCLUDE IN THE MANUSCRIPT.}

This appendix is a separate proof audit. It does not modify the claims or the
numerical tables in \texttt{paper.tex}.
}}
\end{center}

\section{Scope}

This appendix audits two coefficient statements used in the proposed PXL
analysis.  The first statement concerns the degree two Level I shape.  It
shows that a transverse affine restriction has full affine coefficient
support.  The second statement concerns the two degree four Level I shapes.
It gives an exact criterion for full homogeneous coefficient support on a
fixed slice.  The criterion fails after only six variables are fixed.

These are algebraic coefficient statements.  They do not say that the
official coefficient distribution is uniform.  They also do not prove a
solving degree, a Macaulay rank prediction, or a PXL running time.  We record
these limits again after each proof.

\section{The degree two affine coefficient map}
\label{sec:l2}

Let $q=19$ and let $A=\F_{q^2}$.  We regard every $A$ space as an
$\F_q$ space when taking symmetric forms.  Put
\[
  V=A^v,
  \qquad h=v-2.
\]
Let $U$ be the full common column domain and let
$\pi_v:U\to V$ be the public vinegar projection.  Choose an
$\F_q$ subspace $L\subseteq U$ of dimension $h$, with basis
$\ell_1,\ldots,\ell_h$, and define its extension span by
\[
  L_A=\Span_A\{\pi_v(\ell_1),\ldots,\pi_v(\ell_h)\}
  \subseteq V.
\]

\begin{definition}[Extension linear transversality]
\label{def:l2-transverse}
The direction $L$ is extension linearly independent if
\[
  \dim_A L_A=h.
\]
For such a direction, an affine coset $a+L$ is transverse if
\[
  \pi_v(a)\notin L_A.
\]
\end{definition}

The first condition is stronger than injectivity of $\pi_v|_L$ over
$\F_q$.  It says that the $A$ lines generated by the vinegar projections are
independent.  The second condition adds the $A$ line generated by the affine
offset.

\subsection{Uniform coset incidence}

\begin{lemma}[Exact transverse incidence]
\label{lem:l2-incidence}
Fix an extension linearly independent direction $L$ of dimension $h$.
If $a+L$ is a uniform affine coset of $L$ in $U$, then
\[
  \Pr[\pi_v(a)\in L_A]=|A|^{h-v}.
\]
For $h=v-2$, the failure probability is $19^{-4}$.
\end{lemma}

\begin{proof}
A uniform point of $U$ induces the uniform distribution on $U/L$, since
every coset has the same number of representatives.  The event
$\pi_v(a)\in L_A$ is independent of the chosen representative because
$\pi_v(L)\subseteq L_A$.  The induced vinegar coset is therefore uniform in
$V/L_A$.  Exactly one of the $|A|^{v-h}$ cosets is $L_A$ itself.  Hence
\[
  \Pr[\pi_v(a)\in L_A]=\frac{|L_A|}{|V|}=|A|^{h-v}.
\]
Substituting $|A|=19^2$ and $v-h=2$ gives $19^{-4}$.
\end{proof}

This calculation samples a uniform coset first and rejects a nontransverse
coset.  It does not replace the root sampling experiment by a distribution
conditioned on transversality.

\subsection{A blockwise support lemma}

We first isolate the linear algebra used by the coefficient map.  For an
$\F_q$ bilinear form $C:A\times A\to\F_q$, write
\[
  C^{\rm sym}(\xi,\eta)=C(\xi,\eta)+C(\eta,\xi).
\]

\begin{lemma}[The mixed block]
\label{lem:l2-mixed-block}
The map
\[
  \Hom_{\F_q}(A\otimes_{\F_q}A,\F_q)
  \longrightarrow \Sym^2_{\F_q}(A)^*,
  \qquad C\longmapsto C^{\rm sym},
\]
is surjective.  Its kernel is the one dimensional space of alternating
bilinear forms on $A$.
\end{lemma}

\begin{proof}
The characteristic is odd, so a symmetric form $S$ has preimage $S/2$.
The kernel consists exactly of forms satisfying
$C(\xi,\eta)=-C(\eta,\xi)$.  These are the alternating forms.  Since $A$
has dimension two over $\F_q$, their space has dimension
$\binom{2}{2}=1$.
\end{proof}

\begin{proposition}[Full affine support for one degree two base form]
\label{prop:l2-support}
Assume that $L$ is extension linearly independent and that $a+L$ is
transverse.  The restriction map from the fresh native symmetric vinegar
coefficients of one base form to its three descended channel polynomials is
surjective onto
\[
  \F_q[x_1,\ldots,x_h]_{\leq 2}^{\,3}.
\]
Equivalently, every monomial of degree at most two can be assigned an
arbitrary coefficient in $\Sym^2_{\F_q}(A)^*$.
\end{proposition}

\begin{proof}
Put $b=\pi_v(a)$ and $e_i=\pi_v(\ell_i)$.  Transversality gives the direct
sum of $A$ lines
\[
  Ab\oplus Ae_1\oplus\cdots\oplus Ae_h\subseteq V.
\]
Let $B$ be one fresh native symmetric $\F_q$ bilinear form on this sum.  For
$\xi,\eta\in A$, evaluate the corresponding channel on
\[
  b+\sum_{i=1}^h x_i e_i.
\]
The channel label is symmetric in $(\xi,\eta)$, so its coefficient belongs
to $\Sym^2_{\F_q}(A)^*$.  The independent blocks of $B$ supply the
coefficients as follows.

The diagonal block on $Ae_i$ is an arbitrary symmetric form on the two
dimensional $\F_q$ space $A$.  It supplies the coefficient of $x_i^2$.
The block between $Ae_i$ and $Ae_j$ is an arbitrary bilinear form on
$A\times A$.  Its contribution to $x_ix_j$ is its symmetrization.
\Cref{lem:l2-mixed-block} shows that this contribution covers all three
channel coefficients.  Its only unused direction is the one dimensional
alternating kernel.

The block between $Ab$ and $Ae_i$ is governed by the same mixed block map.
It therefore supplies an arbitrary coefficient of $x_i$ in all three
channels.  The diagonal block on $Ab$ supplies the three constant
coefficients.  All listed blocks are independent because the displayed sum
of $A$ lines is direct.  We may set every unused block to zero and extend $B$
arbitrarily to the rest of $V$.  This constructs a preimage for every triple
of affine quadratic polynomials.
\end{proof}

\begin{corollary}[The Level I degree two support statement]
\label{cor:l2-level1}
For $(v,o,\ell,r)=(48,16,2,2)$, take $h=46$.  A uniform coset fails
transversality with probability $19^{-4}$.  Conditional on transversality,
each independent base form maps onto three arbitrary affine quadratics in
$46$ variables.  The sixteen base forms therefore map onto the coefficient
space of $48$ affine quadratics in $46$ variables.
\end{corollary}

\begin{proof}
The probability is \Cref{lem:l2-incidence}.  The native coefficient blocks
for distinct base forms are independent parameters.  Apply
\Cref{prop:l2-support} to each form and take their direct sum.
\end{proof}

\paragraph{Descent to the official coefficient space.}
The block map in \Cref{prop:l2-support} is defined over $\F_q$.  One may
also verify its surjectivity after extending scalars to a splitting field.
If $C$ is its cokernel, scalar extension gives $C\otimes_{\F_q}\Omega=0$.
The extension $\Omega/\F_q$ is faithfully flat, so $C=0$.  Thus the split
channel calculation descends to the official $\F_q$ coefficient constraints.

\paragraph{Distributional limit.}
Surjectivity gives full support when the native coefficient source has full
support.  It gives uniform output only when the native source is uniform on
the relevant finite vector space.  The official byte reduction source is not
uniform on $\F_{19}$.  This proposition does not bound its distance from a
uniform MQ distribution.

\section{The degree four fixed map}
\label{sec:l4}

Put $k=\F_{19}$ and let $\Omega$ be a splitting field for
$A=\F_{19^4}$.  We first work over $\Omega$.  The four split eigenblocks are
spaces $V_a$, with planted oil spaces $O_a\subseteq V_a$, for
$0\leq a<4$.  For both Level I degree four shapes,
\[
  \dim(V_a/O_a)=v=28.
\]
The total block dimensions are $33$ for $(28,5,4,4)$ and $32$ for
$(28,4,4,5)$.

Let $X$ be the space of main variables on a fixed slice.  A realization of
the slice gives four linear maps
\[
  T_a:X\longrightarrow V_a.
\]
Write $q_a:V_a\to V_a/O_a$ for the quotient map and put
\[
  \overline T_a=q_a\circ T_a.
\]

For one native UOV base form, its split coefficient space is the direct sum
of the diagonal blocks
\[
 \mathcal B_{aa}
 =\{B\in\Sym^2(V_a^*):B|_{O_a\times O_a}=0\}
\]
and the cross blocks
\[
 \mathcal B_{ab}
 =\{C\in\Hom(V_a\otimes V_b,\Omega):
       C|_{O_a\times O_b}=0\},
 \qquad a<b.
\]
Put
\[
 \mathcal B
 =\bigoplus_a\mathcal B_{aa}
  \oplus\bigoplus_{a<b}\mathcal B_{ab}.
\]
Equivalently, put $V=\bigoplus_aV_a$ and $O=\bigoplus_aO_a$.  Then
$\mathcal B$ is the space of symmetric forms on $V$ that vanish on
$O\times O$, written in its four diagonal and six cross blocks.
The fixed realization $T=(T_0,T_1,T_2,T_3)$ defines the linear coefficient
map
\[
 \Phi_T:\mathcal B
 \longrightarrow
 \bigoplus_{0\leq a\leq b<4}\Sym^2(X^*).
 \label{eq:l4-fixed-map}
\]
On a diagonal block it is pullback by $T_a$.  On a cross block its value is
the symmetric form
\[
 (x,y)\longmapsto
 C(T_a x,T_b y)+C(T_a y,T_b x).
 \label{eq:l4-cross-map}
\]
The target has four diagonal channels and six cross channels.

\subsection{The fixed map criterion}

\begin{lemma}[One diagonal block]
\label{lem:l4-fixed-diagonal}
Let $T:X\to V$ and let $O\subseteq V$.  The pullback map
\[
 \{B\in\Sym^2(V^*):B|_{O\times O}=0\}
 \longrightarrow\Sym^2(X^*)
\]
is onto if and only if the composite $X\to V/O$ is injective.
\end{lemma}

\begin{proof}
Suppose that a nonzero $x\in X$ maps into $O$.  Every form in the image then
vanishes at $(x,x)$.  This is a proper linear condition on
$\Sym^2(X^*)$, so the pullback map is not onto.

Conversely, suppose that $X\to V/O$ is injective.  Choose a complement
$V=O\oplus C$.  The projection of $T(X)$ to $C$ is an isomorphic copy of
$X$.  Given $S\in\Sym^2(X^*)$, define a symmetric form on this copy by
transporting $S$.  Extend it to $C$, and set its blocks involving $O$ to
zero.  The resulting form vanishes on $O\times O$ and pulls back to $S$.
\end{proof}

\begin{lemma}[One cross block]
\label{lem:l4-fixed-cross}
Let $T_a:X\to V_a$ and $T_b:X\to V_b$.  If both maps
$X\to V_a/O_a$ and $X\to V_b/O_b$ are injective, then the cross block map
in \eqref{eq:l4-cross-map} is onto $\Sym^2(X^*)$.
\end{lemma}

\begin{proof}
Choose complements $V_a=O_a\oplus C_a$ and
$V_b=O_b\oplus C_b$.  Let $c_a:X\to C_a$ and $c_b:X\to C_b$ be the
projections of $T_a$ and $T_b$.  Both maps are injective.  Given
$S\in\Sym^2(X^*)$, define a bilinear form on
$c_a(X)\times c_b(X)$ by
\[
 C(c_a x,c_b y)=\frac{1}{2}S(x,y).
\]
Extend this form to $C_a\times C_b$, and set every block involving an oil
space to zero.  The extension lies in $\mathcal B_{ab}$.  Its symmetric
pullback is $S$ because the characteristic is odd.
\end{proof}

\begin{theorem}[Fixed map linear surjectivity]
\label{thm:l4-fixed-surjectivity}
The map $\Phi_T$ in \eqref{eq:l4-fixed-map} is onto if and only if every
quotient map
\[
  \overline T_a:X\longrightarrow V_a/O_a
\]
is injective.
\end{theorem}

\begin{proof}
The source and target of $\Phi_T$ are direct sums over the ten channel
blocks.  If every $\overline T_a$ is injective, the four diagonal blocks are
onto by \Cref{lem:l4-fixed-diagonal}.  The six cross blocks are onto by
\Cref{lem:l4-fixed-cross}.  The blocks use independent source coordinates,
so their direct sum is onto.

For the converse, suppose that some $\overline T_a$ is not injective.  The
projection of $\Phi_T$ to diagonal channel $(a,a)$ is not onto by
\Cref{lem:l4-fixed-diagonal}.  Hence $\Phi_T$ is not onto.
\end{proof}

For five independent native base forms, the full map to the $50$ channel
quadrics is the direct sum of five copies of $\Phi_T$.  It is onto under the
same four injectivity conditions.

\begin{corollary}[The six variable fixing cannot give full support]
\label{cor:l4-six-fixing-obstruction}
For both Level I degree four shapes, fixing six variables in a
$37$ dimensional slice leaves $\dim X=31$.  The fixed coefficient map
$\Phi_T$ is not onto for any realization $T$.
\end{corollary}

\begin{proof}
Every quotient $V_a/O_a$ has dimension $28$.  No map from a
$31$ dimensional space to a $28$ dimensional space is injective.  Apply
\Cref{thm:l4-fixed-surjectivity}.
\end{proof}

Thus six fixed variables do not give full homogeneous coefficient support.
This criterion allows full support only after the main variable space has
dimension at most $28$.  Starting from dimension $37$, this requires fixing
at least nine variables and choosing a realization that satisfies the four
injectivity conditions.

\subsection{The open good locus and realization}

\begin{proposition}[Open nonempty locus]
\label{prop:l4-open-locus}
Let $n=\dim X$ and let
\[
 \mathcal P=\prod_{a=0}^3\Hom(X,V_a)
\]
be the space of abstract four block realizations.  The locus on which
$\Phi_T$ is onto is Zariski open.  It is nonempty in $\mathcal P$ if and only
if $n\leq v$.

If the SNOVA construction supplies only a realization locus
$\mathcal R\subseteq\mathcal P$, the good locus in that family is
$\mathcal R\cap\mathcal P^{\rm good}$.  It is open in $\mathcal R$.  Its
nonemptiness requires one realizable tuple with all four quotient maps
injective.
\end{proposition}

\begin{proof}
By \Cref{thm:l4-fixed-surjectivity}, the good locus is the intersection of
the four full rank loci for $q_aT_a$.  Each full rank condition is Zariski
open.  If $n>v$, every one of these loci is empty.  If $n\leq v$, choose an
injection $X\to V_a/O_a$ for each $a$ and lift it through $q_a$.  This gives
a point in the good locus of $\mathcal P$.

The last statement follows by intersecting with $\mathcal R$.  Nonemptiness
in the ambient product does not show that the constrained realization family
meets the good locus.
\end{proof}

\subsection{Descent to the base field}

\begin{proposition}[Faithfully flat descent for a realized map]
\label{prop:l4-descent}
Suppose there are finite dimensional $k$ spaces $\mathcal B_k$ and
$\mathcal C_k$, and a $k$ linear map
\[
 \Phi_{T,k}:\mathcal B_k\longrightarrow\mathcal C_k,
\]
whose scalar extension to $\Omega$ identifies with the assembled split map
$\Phi_T$.  Then $\Phi_{T,k}$ is onto if and only if $\Phi_T$ is onto.
\end{proposition}

\begin{proof}
Let $C$ be the cokernel of $\Phi_{T,k}$.  Scalar extension is exact, so the
cokernel of $\Phi_T$ is $C\otimes_k\Omega$.  If $\Phi_T$ is onto, this tensor
product is zero.  The extension $\Omega/k$ is faithfully flat, so $C=0$.
The converse follows from right exactness of scalar extension.
\end{proof}

The qualification in \Cref{prop:l4-descent} is necessary.  The individual
eigenblocks and oil spaces can exist only after passing to the splitting
field.  One must first show that the assembled ten channel map is the scalar
extension of a map defined over $k$.  An abstract good tuple over $\Omega$
does not produce a $k$ rational slice, and faithful flatness does not prove
that the official byte reduction source is uniform.

\subsection{Varying slice incidence calculations}
\label{sec:l4-varying-slice}

The rest of this section records the earlier codimension and tangent
calculations.  These calculations allow the slice intersections to vary with
the target quadrics.  They do not prove that the coefficient map for a fixed
slice is onto, and they do not change
\Cref{cor:l4-six-fixing-obstruction}.

Let $Z$ be the $37$ dimensional slice before the six PXL coefficient
variables are fixed.  Write
\[
  Z=E\oplus X,
  \qquad \dim E=6,
  \qquad \dim X=31.
\]
The four Frobenius eigenprojections are maps
\[
  \rho_a:Z\longrightarrow V_a,
  \qquad 0\leq a<4.
\]
For the shape $(28,5,4,4)$, a generic $\rho_a$ has a four dimensional
kernel.  For the shape $(28,4,4,5)$, it has a five dimensional kernel.

\begin{lemma}[Diagonal channel image on the varying slice]
\label{lem:l4-diagonal-kernel}
Let $\rho:Z\to V$ be surjective with kernel $K$.  Pullback identifies
$\Sym^2(V^*)$ with the symmetric forms $S$ on $Z$ for which
$K\subseteq\rad S$.  If $\dim Z=37$ and $\dim K=c$, its codimension in
$\Sym^2(Z^*)$ is
\[
  \binom{38}{2}-\binom{38-c}{2}.
\]
This codimension is $142$ for $c=4$ and $175$ for $c=5$.
\end{lemma}

\begin{proof}
A pulled back form vanishes whenever either input lies in $K$.  Conversely,
such a form descends uniquely to $Z/K\cong V$.  The dimension is therefore
$\binom{38-c}{2}$.  Subtraction from
$\dim\Sym^2(Z^*)=\binom{38}{2}$ gives the formula and the two stated values.
\end{proof}

\begin{lemma}[Alternating cross lift on the varying slice]
\label{lem:l4-cross-lift}
Let $K_a,K_b\subseteq Z$ be disjoint subspaces of dimension $c$.  A
symmetric form $S$ on $Z$ is induced by a bilinear block that descends
through
\[
  (Z/K_a)\times(Z/K_b)
\]
if and only if
\[
  S|_{K_a\times K_a}=0,
  \qquad S|_{K_b\times K_b}=0.
\]
The cross channel image has codimension
\[
  2\binom{c+1}{2}=c(c+1).
\]
This codimension is $20$ for $c=4$ and $30$ for $c=5$.
\end{lemma}

\begin{proof}
The two vanishing conditions are necessary.  For sufficiency, choose an
alternating form $A_0$ whose values on $K_a\times Z$ and $Z\times K_b$
cancel the corresponding values of $S$.  The prescriptions agree on
$K_a\times K_b$.  They are compatible with alternation on $K_a$ and $K_b$
because the two displayed restrictions of $S$ vanish.  Since
$K_a\cap K_b=0$, extend this partial alternating form to all of $Z$.

The bilinear form $C=S+A_0$ vanishes on $K_a\times Z$ and on
$Z\times K_b$.  It therefore descends through the two quotients.  Its
quadratic polynomial is unchanged because $A_0(x,x)=0$.  The restriction
maps to $\Sym^2(K_a^*)$ and $\Sym^2(K_b^*)$ are independent for disjoint
subspaces.  Their total dimension gives the codimension.
\end{proof}

There are four diagonal channels and six cross channels for each base form.
The preceding lemmas give the following codimension audit.

\begin{table}[h]
\centering
\caption{Codimension of the split coefficient family on the $37$ dimensional
varying slice.}
\label{tab:l4-codim}
\begin{tabular}{@{}lrrrr@{}}
\toprule
Shape & Diagonal & Cross & Per base form & Five base forms\\
\midrule
$(28,5,4,4)$ & $142$ & $20$ & $4(142)+6(20)=688$ & $3440$\\
$(28,4,4,5)$ & $175$ & $30$ & $4(175)+6(30)=880$ & $4400$\\
\bottomrule
\end{tabular}
\end{table}

\begin{lemma}[Injectivity before the oil quotient]
\label{lem:l4-main-injective}
There is a nonempty Zariski open set of decompositions $Z=E\oplus X$ with
$\dim X=31$ for which
\[
  X\cap\ker\rho_a=0
  \qquad\text{for every }a.
\]
This holds for both Level I degree four shapes.
\end{lemma}

\begin{proof}
For one kernel of dimension $c$, the condition $X\cap K=0$ is open on
$\Gr(31,Z)$.  It is nonempty when $31+c\leq37$.  Here $c$ is four or five.
Each such open set is dense because the Grassmannian is irreducible.  Their
intersection is therefore nonempty.  On this open set, every
$\rho_a|_X$ is injective.
\end{proof}

Let $O_a\subseteq V_a$ have dimension five for $(28,5,4,4)$ and four for
$(28,4,4,5)$.  For a generic injective image $\rho_a(X)$, its intersection
with $O_a$ has dimension
\[
 \dim(\rho_a(X)\cap O_a)=31+\dim O_a-\dim V_a=3.
\]
Its pullback is
\[
 I_a=\ker(X\xrightarrow{\rho_a}V_a\to V_a/O_a).
\]
Thus injectivity before the oil quotient is compatible with a three
dimensional kernel after that quotient.  Every one of the five diagonal base
forms vanishes on the fixed space $I_a$.

The next incidence calculation permits this three space to vary with the
five target quadrics.

\begin{proposition}[Five quadrics with a varying common three space]
\label{prop:l4-incidence}
Let $X$ have dimension $31$ over an algebraically closed field of odd
characteristic.  The projection
\[
 \left\{(I,Q_1,\ldots,Q_5):
 I\in\Gr(3,X),\ Q_s|_{I\times I}=0\right\}
 \longrightarrow \Sym^2(X^*)^5
\]
is dominant.
\end{proposition}

\begin{proof}
The Grassmannian has dimension
\[
  \dim\Gr(3,31)=3(31-3)=84.
\]
For a fixed $I$, the five isotropy conditions impose
$5\binom{4}{2}=30$ linear equations.  This count leaves room for dominance,
but it does not prove it.  We now exhibit a point where the projection has
surjective differential.

Choose a three space $I$.  In $X/I$, choose five pairwise disjoint three
spaces $W_1,\ldots,W_5$.  Their total dimension is $15$, while
$\dim(X/I)=28$.  For each $s$, choose $Q_s$ so that its pairing between $I$
and $W_s$ is perfect.  Set its pairing between $I$ and $W_t$ to zero when
$t\ne s$, and impose $Q_s|_{I\times I}=0$.

A tangent vector to the Grassmannian at $I$ is a map
$\delta:I\to X/I$.  The derivative of the $s$th isotropy equation is
\[
  (e,e')\longmapsto
  Q_s(\delta e,e')+Q_s(e,\delta e').
\]
The component of $\delta$ in $W_s$ can be chosen independently of its
components in the other $W_t$.  The perfect pairing identifies
$\Hom(I,W_s)$ with $\Hom(I,I^*)$.  Under this identification, the displayed
derivative is symmetrization.  Symmetrization is onto $\Sym^2(I^*)$ because
two is invertible.  Thus the derivative from
$\Hom(I,X/I)$ onto $\Sym^2(I^*)^5$ is surjective.

Given any first order change in the five quadrics, choose $\delta$ to cancel
their restrictions to $I$.  The incidence projection is smooth and
submersive at the constructed point.  Its image contains a nonempty open set,
which proves dominance.
\end{proof}

The same tangent construction can impose pairwise disjoint varying spaces
$I_0,I_1,I_2,I_3$.  Pairwise disjointness is open, and four three spaces fit
in $X$.  This remains a statement about the incidence projection in which the
four spaces vary with the target tuple.

\begin{lemma}[Oil compatible cross extension for varying spaces]
\label{lem:l4-oil-cross}
Let $I_a,I_b\subseteq X$ be disjoint three spaces.  Every symmetric form
$S$ on $X$ is the quadratic form of a bilinear block $C$ satisfying
\[
  C|_{I_a\times I_b}=0.
\]
\end{lemma}

\begin{proof}
Define an alternating form $A_0$ on $I_a\oplus I_b$ by
\[
  A_0(u,v)=-S(u,v)
  \qquad (u\in I_a,\ v\in I_b),
\]
and extend it to $X$.  Put $C=S+A_0$.  Then $C$ vanishes on
$I_a\times I_b$.  Since $A_0(x,x)=0$, the quadratic polynomial of $C$ is the
one defined by $S$.
\end{proof}

\begin{remark}[Why the varying slice calculation is not a fixed map theorem]
For a fixed realization $T$, the spaces $I_a=\ker\overline T_a$ are fixed.
The diagonal image then consists of forms that vanish on these fixed spaces.
The incidence projection in \Cref{prop:l4-incidence} instead chooses $I_a$
after seeing the target five tuple.  It can therefore be dominant even though
every fixed fiber is a proper linear subspace.  The cross extension in
\Cref{lem:l4-oil-cross} has the same varying space qualification.  These
calculations may support an analysis that varies the slice with the
coefficients, but they do not prove surjectivity or dominance for the fixed
map $\Phi_T$.
\end{remark}

\section{The Hilbert coefficient calculation}
\label{sec:hilbert}

The homogeneous model has $50$ quadrics in $31$ variables.  Its formal
semi regular Hilbert series is
\[
  H_{31}(t)=\frac{(1-t^2)^{50}}{(1-t)^{31}}.
\]
Expanding through degree eight gives
\begin{align*}
 H_{31}(t)={}&1+31t+446t^2+3906t^3+22801t^4
 +89807t^5\\
 &+216992t^6+139872t^7-1166808t^8+O(t^9).
\end{align*}
The first negative coefficient occurs at degree eight.

If one homogenizes the affine system with one extra variable, the formal
series becomes
\[
  H_{32}(t)=\frac{(1-t^2)^{50}}{(1-t)^{32}}.
\]
Its expansion is
\begin{align*}
 H_{32}(t)={}&1+32t+478t^2+4384t^3+27185t^4
 +116992t^5\\
 &+333984t^6+473856t^7-692952t^8+O(t^9).
\end{align*}
This model also selects degree eight.

\begin{proposition}[Exact coefficient formula]
\label{prop:hilbert-coeff}
For $n\in\{31,32\}$, the coefficient of $t^d$ in $H_n(t)$ is
\[
  [t^d]H_n(t)
  =\sum_{j=0}^{\lfloor d/2\rfloor}
    (-1)^j\binom{50}{j}\binom{n+d-2j-1}{d-2j}.
\]
The two displayed expansions follow from this formula.
\end{proposition}

\begin{proof}
Use the binomial expansions
\[
 (1-t^2)^{50}=\sum_{j=0}^{50}(-1)^j\binom{50}{j}t^{2j},
 \qquad
 (1-t)^{-n}=\sum_{r\geq0}\binom{n+r-1}{r}t^r,
\]
and collect the terms with $2j+r=d$.
\end{proof}

\section{Logical status}

The results in this appendix have the following scope.

\begin{enumerate}
\item The degree two theorem proves exact affine coefficient support on every
      transverse coset.  It also gives the exact $19^{-4}$ incidence loss for
      the Level I direction of dimension $46$.

\item The degree four theorem gives an exact fixed map criterion.  The ten
      channel map is onto if and only if all four maps to the vinegar
      quotients $V_a/O_a$ are injective.

\item After six variables are fixed, the main space has dimension $31$ and
      each vinegar quotient has dimension $28$.  Full homogeneous coefficient
      support is therefore impossible for every fixed realization.  Fixing at
      least nine variables removes this dimension obstruction, but one must
      still prove that the constrained realization locus meets the open good
      locus.

\item The codimension and tangent incidence calculations in
      \Cref{sec:l4-varying-slice} permit the slice intersections to vary with
      the target quadrics.  They do not prove dominance for a fixed
      coefficient map.

\item Faithfully flat descent applies only after the assembled split map has
      been identified as the scalar extension of a base field map.  It does
      not produce a base field slice and does not give a distributional claim
      for the official finite source.

\item The Hilbert expansions are exact formal calculations.  The claim that a
      production system follows either series through degree eight remains the
      ordinary semi regularity assumption.  The fixed map theorem does not
      imply it.

\item This appendix does not prove the Macaulay matrix ranks used by PXL.  It
      does not prove that symbolic preprocessing and later variable fixing
      preserve the random MQ rank model.  It does not provide a PXL operation
      count.

\item Surjectivity gives full support only when the native source has full
      support.  Uniform output further requires a uniform native source.  The
      official byte reduction source is not uniform on $\F_{19}$.
\end{enumerate}

These limits separate the algebraic proof from the solver model.  Any later
manuscript statement must keep the same separation.

\end{document}
